% ==STYLE-BEGIN==
\documentclass[UTF8,11pt]{ctexart}
\usepackage{amsmath,amssymb,amsthm}
\usepackage[a4paper,margin=2.4cm,top=2.3cm,bottom=2.3cm]{geometry}
\usepackage{fancyhdr}
\usepackage{enumitem}
\usepackage{hyperref}
\hypersetup{hidelinks}

% ---- 版式 ----
\linespread{1.15}
\setlength{\parskip}{0.35em}
\setlength{\parindent}{1.6em}
\setlist{nosep,leftmargin=2em}

% ---- 页眉 ----
\pagestyle{fancy}
\fancyhf{}
\fancyhead[L]{\small\coursename}
\fancyhead[R]{\small\thepage}
\renewcommand{\headrulewidth}{0.4pt}

% ---- 定理环境：中文用正体，标题加粗，按小节共用一条流水号 ----
\newtheoremstyle{cnthm}{0.7em}{0.7em}{\normalfont}{}{\bfseries}{}{0.5em}{}
\theoremstyle{cnthm}
\newtheorem{thm}{定理}[section]
\newtheorem{lem}[thm]{引理}
\newtheorem{prop}[thm]{命题}
\newtheorem{cor}[thm]{推论}
\newtheorem{defn}[thm]{定义}
\newtheorem{rem}[thm]{注}
\newtheorem{exa}[thm]{例}
\newtheorem{conj}[thm]{猜想}
\renewcommand{\proofname}{\bfseries 证明}
\providecommand{\coursename}{}
% ==STYLE-END==
\renewcommand{\coursename}{解析几何\,·\,0916}
\title{解析几何 0916}
\date{}
\begin{document}
\maketitle

\section{\ Recall：线性组合、线性表示、线性相关与线性无关}

\begin{itemize}
\item \textbf{共线}\quad $\vec a$ 与 $\vec b$ 共线
$\Longleftrightarrow$ $\exists$ 不全为零的实数
$\lambda,\mu$，s.t.
$\lambda\vec a+\mu\vec b=\vec 0$；
\item \textbf{共面}\quad $\vec a,\vec b,\vec c$ 共面
$\Longleftrightarrow$ $\exists$ 不全为 $0$ 的数
$\lambda,\mu,\nu$，s.t.
$\lambda\vec a+\mu\vec b+\nu\vec c=\vec 0$；
\item \textbf{线性相关}\quad
$\vec a_1,\vec a_2,\ldots,\vec a_n$，若存在不全为零的实数
$k_1,k_2,\ldots,k_n$，s.t.
$k_1\vec a_1+k_2\vec a_2+\cdots+k_n\vec a_n=\vec 0$，
则称 $\vec a_1,\ldots,\vec a_n$ 线性相关；
\item \textbf{线性无关}\quad
$\vec a_1,\vec a_2,\ldots,\vec a_n$，若不存在不全为零的实数
$k_1,k_2,\ldots,k_n$，s.t.
$k_1\vec a_1+k_2\vec a_2+\cdots+k_n\vec a_n=\vec 0$，
则称 $\vec a_1,\ldots,\vec a_n$ 线性无关；
\item 向量组 $\vec a_1,\vec a_2,\ldots,\vec a_n$ 线性相关／无关
$\Longleftrightarrow$ 向量方程
$k_1\vec a_1+k_2\vec a_2+\cdots+k_n\vec a_n=0$
有／无非零解，or 方程解不唯一／唯一；
\item $\vec a_1,\vec a_2,\ldots,\vec a_n$ 线性相关
$\Longleftrightarrow$ $\exists$ 某个 $\vec a_i$，s.t.
$\vec a_i$ 能被其余向量线性表示（有多余向量）。
\end{itemize}

\section{\ 向量的坐标表示}

\begin{thm}
若 $\vec a,\vec b,\vec c$ 不共面（线性无关），
则 $\forall\vec d$，存在唯一的实数组 $\lambda,\mu,\nu$，s.t.
\[
\vec d=\lambda\vec a+\mu\vec b+\nu\vec c .
\]
\end{thm}

在空间直角坐标系 $[O,X,Y,Z]$ 中，取 $OX,OY,OZ$ 轴方向的单位向量
$\vec e_1,\vec e_2,\vec e_3$，则 $\forall\vec a$，存在唯一的
$x,y,z$，s.t.
\[
\vec a=x\vec e_1+y\vec e_2+z\vec e_3 .
\]
称 $(x,y,z)$ 是 $\vec a$ 在
$[O;\vec e_1,\vec e_2,\vec e_3]$ 的坐标，记作
$\vec a=(x,y,z)$。

若 $\vec a=(x_1,y_1,z_1)$，$\vec b=(x_2,y_2,z_2)$，则有
\[
\begin{aligned}
\vec a+\vec b
&=(x_1\vec e_1+y_1\vec e_2+z_1\vec e_3)
+(x_2\vec e_1+y_2\vec e_2+z_2\vec e_3)\\
&=(x_1\vec e_1+x_2\vec e_1)+(y_1\vec e_2+y_2\vec e_2)
+(z_1\vec e_3+z_2\vec e_3)\\
&=(x_1+x_2)\vec e_1+(y_1+y_2)\vec e_2+(z_1+z_2)\vec e_3 ,
\end{aligned}
\]
于是
$\vec a+\vec b=(x_1+x_2,\ y_1+y_2,\ z_1+z_2)$。
同样地
\[
\begin{aligned}
\lambda\vec a
&=\lambda(x_1\vec e_1+y_1\vec e_2+z_1\vec e_3)\\
&=(\lambda x_1)\vec e_1+(\lambda y_1)\vec e_2+(\lambda z_1)\vec e_3\\
\Longrightarrow\quad
\lambda\vec a&=(\lambda x_1,\ \lambda y_1,\ \lambda z_1).
\end{aligned}
\]

$\vec e_1,\vec e_2,\vec e_3$ 称为基向量。

\section{\ 向量的内积（点乘）前置}

\subsubsection*{投影}

若 $\vec e$ 是单位向量，$\vec a$ 是空间向量，此时
$\vec e,\vec a$ 可以确定一个平面 $\pi$，在 $\pi$ 中取与
$\vec e$ 垂直的单位向量 $\vec e'$，于是
\[
\vec a=\lambda\vec e+\mu\vec e'
=\vec a_{\vec e}+\vec a_h .
\]

\begin{itemize}
\item 称 $\vec a_{\vec e}$ 是 $\vec a$ 在 $\vec e$ 方向上的投影
（内支量，水平分量）；
\item 称 $\vec a_h$ 是 $\vec a$ 在 $\vec e$ 方向上的外支量
（垂直分量）；
\item 称 $\lambda$ 是投影值。
\end{itemize}

从映射的角度看投影：
\[
P_{\vec e}:V\to V,\qquad \vec a\longmapsto\vec a_{\vec e} .
\]

\textbf{性质}
\[
\begin{aligned}
&1.\quad P_{\vec e}(\vec a+\vec b)
=P_{\vec e}(\vec a)+P_{\vec e}(\vec b);\\
&2.\quad P_{\vec e}(\lambda\vec a)
=\lambda P_{\vec e}(\vec a).
\end{aligned}
\]
证明将向量拆成水平分量和垂直分量即可，不在此赘述。

由此可见，$f:V\to V$ 满足
\[
\begin{aligned}
&1.\quad f(\vec a+\vec b)=f(\vec a)+f(\vec b);\\
&2.\quad f(\lambda\vec a)=\lambda f(\vec a),
\end{aligned}
\]
称 $f$ 是\textbf{线性映射}（保持加法与数乘）。

\subsubsection*{夹角}

$\vec a$ 与 $\vec b$ 的夹角定义成 $\vec a,\vec b$ 方向的夹角，
用 $\langle\vec a,\vec b\rangle$ 表示。

\textbf{性质}
\[
\begin{aligned}
&1.\quad 0\leq\langle\vec a,\vec b\rangle\leq\pi;\\
&2.\quad \langle\vec a,\vec b\rangle
=\langle\vec b,\vec a\rangle;\\
&3.\quad \langle\lambda\vec a,\vec b\rangle=
\begin{cases}
\langle\vec a,\vec b\rangle, & \lambda>0,\\
\pi-\langle\vec a,\vec b\rangle, & \lambda<0 .
\end{cases}
\end{aligned}
\]

\section{\ 向量的内积}

$\vec a$ 与 $\vec b$ 的内积 $\vec a\cdot\vec b$ 定义为
$|\vec a||\vec b|\cos\langle\vec a,\vec b\rangle$。

$\vec a\cdot\vec e
=|\vec a||\vec e|\cos\langle\vec a,\vec e\rangle
=|\vec a|\cos\langle\vec a,\vec e\rangle
=\vec a$ 在 $\vec e$ 方向的投影值。

于是有
\[
\vec a_{\vec e}=(\vec a\cdot\vec e)\vec e .
\]
注意到
\[
\vec a_{\vec e}\cdot\vec e
=\bigl((\vec a\cdot\vec e)\vec e\bigr)\cdot\vec e
=\vec a\cdot\vec e .
\]

\subsubsection*{内积的性质}

\[
\begin{aligned}
&1.\quad \vec a\cdot\vec b=\vec b\cdot\vec a
&&\text{交换律};\\
&2.\quad (\lambda\vec a)\cdot\vec b
=\lambda(\vec a\cdot\vec b) &&\text{结合律};\\
&3.\quad (\vec a+\vec b)\cdot\vec e
=\vec a\cdot\vec e+\vec b\cdot\vec e &&\text{分配律}.
\end{aligned}
\]
1、2 证明完全展开即可，在此不再赘述；3 证明见下：
\[
\begin{aligned}
(\vec a+\vec b)\cdot\vec e
&=(\vec a+\vec b)_{\vec e}\cdot\vec e\\
&=(\vec a_{\vec e}+\vec b_{\vec e})\cdot\vec e\\
&=\bigl((\vec a\cdot\vec e)\vec e
+(\vec b\cdot\vec e)\vec e\bigr)\cdot\vec e\\
&=\bigl((\vec a\cdot\vec e+\vec b\cdot\vec e)\vec e\bigr)\cdot\vec e\\
&=(\vec a\cdot\vec e+\vec b\cdot\vec e)(\vec e\cdot\vec e)\\
&=\vec a\cdot\vec e+\vec b\cdot\vec e .
\end{aligned}
\]

由此可见 $f:V\times V\to\mathbb R$ 满足
\[
\begin{aligned}
&1.\quad f(\vec a,\vec b)=f(\vec b,\vec a)
&&\text{对称性};\\
&2.\quad f(\lambda\vec a,\vec b)=\lambda f(\vec a,\vec b)
&&\text{双线性性 1};\\
&3.\quad f(\vec a_1+\vec a_2,\vec b)
=f(\vec a_1,\vec b)+f(\vec a_2,\vec b)
&&\text{双线性性 2};\\
&4.\quad \forall\vec a\neq\vec 0,\ f(\vec a,\vec a)>0
&&\text{正定性}.
\end{aligned}
\]

在 $[O;\vec e_1,\vec e_2,\vec e_3]$ 中，若
$\vec a=(x_1,y_1,z_1)$，$\vec b=(x_2,y_2,z_2)$，有
\[
\begin{aligned}
\vec a\cdot\vec b
&=(x_1\vec e_1+y_1\vec e_2+z_1\vec e_3)
(x_2\vec e_1+y_2\vec e_2+z_2\vec e_3)\\
&=x_1x_2+y_1y_2+z_1z_2 .
\end{aligned}
\]

克氏符号：
\[
\delta_{ij}=
\begin{cases}
1, & i=j,\\
0, & i\neq j .
\end{cases}
\]
用克氏符号表示向量点乘的坐标形式：
\[
\begin{aligned}
\vec a\cdot\vec b
&=\left(\sum_{i=1}^{3}x_i\vec e_i\right)
\left(\sum_{j=1}^{3}y_j\vec e_j\right)
=\sum_{i,j=1}^{3}x_iy_j\vec e_i\vec e_j\\
&=\sum_{i,j=1}^{3}\delta_{ij}x_iy_j
=\sum_{i=1}^{3}x_iy_i .
\end{aligned}
\]

\end{document}
